\documentclass[11pt]{article}
% Every number and every result sentence comes from macros.tex, generated by `python -m nycemis paper`
% from results files. Do not type results into this file.
\usepackage[margin=1in]{geometry}
\usepackage{graphicx}
\usepackage{booktabs}
\usepackage{amsmath}
\usepackage[hidelinks]{hyperref}
\usepackage{natbib}
\input{macros.tex}
\input{macros_sensitivity.tex}

\title{Does forecast-driven eco-rerouting reduce simulated traffic emissions?\\
\large An audited, preregistered SUMO study of Flushing, Queens}
\author{Mohammed Y. Hossain}
\date{}

\begin{document}
\maketitle

\begin{abstract}
A 2024 student project reported that machine-learning-triggered rerouting reduced simulated NOx in Flushing,
Queens, by \LegacyNOxDropRF\% (random forest) and \LegacyNOxDropXGB\% (XGBoost). An audit of its notebook shows
that the comparison set a run of about \LegacyBaselineHorizon~s against runs stopped after
\LegacyControlledSteps~s, and that repeated rerouting calls returned unchanged routes. We rebuild the study under a
matched, preregistered protocol in SUMO: the same network, demand, fleet and seed in every condition, runs to
completion, paired seeds, a travel-time rerouting baseline and a random-rerouting placebo. Replicating the legacy
rule under this protocol, the legacy metric still reports a large drop while the matched metric does not, and the
rule almost never changes a route. We then test emission-aware rerouting driven by a \HorizonS-second traffic-state
forecast against preregistered comparisons at three demand levels. Against no rerouting and against travel-time
rerouting it lowered simulated total NOx at every level (Section~\ref{sec:results}), \CompositionWord\ by moving
traffic onto residential and service streets, whose NOx changed by \ResidNOxLight\%, \ResidNOxModerate\% and
\ResidNOxHeavy\% (light, moderate, heavy demand) in a network whose speed coding favours such detours. On held-out
seeds the forecaster \ForecastVerdictR\ a per-edge historical average of the emission rate it supplies
(Section~\ref{sec:forecast}). On a variant of the network with New York City's default 25~mph limit and stop control
at minor-street junctions, measured-emission rerouting no longer lowered NOx against travel-time rerouting but raised
it, by \SensEcoReactiveVsTtAdaptiveTotalNoxGLightPosted\%, \SensEcoReactiveVsTtAdaptiveTotalNoxGModeratePosted\% and
\SensEcoReactiveVsTtAdaptiveTotalNoxGHeavyPosted\% (Section~\ref{sec:sensitivity}): the simulated benefit depends on
how the network is coded. All
quantities are simulated tailpipe emissions, not air quality or exposure.
\end{abstract}

\section{Introduction}
Rerouting vehicles away from congested or high-emitting streets is a natural lever for lowering urban traffic
emissions, and microscopic simulation makes such interventions cheap to test. It also makes it easy to measure
them wrongly. This paper starts from an audit of a student project whose headline result could not survive two
questions: were the compared runs the same length, and did the intervention change anything? We answer both with
the original notebook's saved outputs and with a replication, then ask the substantive question with a protocol
that removes those failure modes: does emission-aware rerouting, driven by a forecast of the next minute of
traffic, reduce network NOx compared with no rerouting, with travel-time rerouting, and with measurement-driven
eco-rerouting, and at what cost in travel time?

\section{Related work}
Route choice changes vehicle energy use and emissions, and the fastest route is not always the lowest-emission one
\citep{ahn2008route}. Eco-routing has been evaluated network-wide in simulation \citep{ahn2013ecorouting}.
Learning-based control has been studied for traffic signals with attention to air quality
\citep{fazzini2022signals,haydari2021signals} and for dynamic rerouting aimed at travel efficiency
\citep{du2021rerouting}. We use SUMO \citep{lopez2018sumo} and LightGBM \citep{ke2017lightgbm}; the statistics
use the Wilcoxon signed-rank test \citep{wilcoxon1945}, two one-sided tests for equivalence
\citep{schuirmann1987tost}, BCa bootstrap intervals \citep{efron1987bca} and Holm's procedure \citep{holm1979}.

\section{Methods}
\subsection{Scenario, demand and fleet}
The network is built from OpenStreetMap for downtown Flushing with the options of SUMO's OSM Web Wizard, clipped to
the study box, with two documented repairs (a \JunctionJoinM~m junction-join distance and SUMO's junction-blocker option at \JunctionBlockerS~s) that
remove deadlocks at one junction cluster. Where OpenStreetMap has no speed tag, netconvert assigns SUMO's default
speed for the road type, as the Web Wizard does: \SpeedSharePrimSec\% of primary and secondary lane-km is coded at
\SpeedPrimSec~km/h, \SpeedShareTertiary\% of tertiary lane-km at \SpeedTertiary~km/h and \SpeedShareResid\% of
residential lane-km at \SpeedResid~km/h. The network has \JunctionsSignal\ signalised, \JunctionsPriority\ priority
and \JunctionsStop\ stop-controlled junctions; Section~\ref{sec:discussion} returns to what this implies. Demand is generated per vehicle class with the Web Wizard's trip
options and routed explicitly; three levels are calibrated on pilot seeds to a relative speed ratio (ideal over
actual travel time, relative to a near-empty network) of \RsrLight, \RsrModerate\ and \RsrHeavy\ (light,
moderate, heavy), i.e.\ \DensityLight, \DensityModerate\ and \DensityHeavy\ vehicles per hour per km of road.
Emission classes come from HBEFA~4.2 as implemented in SUMO; each vehicle's class is fixed in the route file, so all
conditions share the fleet.

\subsection{Matched protocol}
Every condition of a seed shares the network, the byte-identical route file, the SUMO seed and options; runs
continue until every vehicle has arrived. A run is valid if at least \ValidCompletionPct\% of vehicles complete and at most
\ValidTeleportPct\% teleport. The primary outcome is total NOx from SUMO's per-trip emissions.

\subsection{Controllers}
C0 keeps the free-flow routes. C1 is SUMO's travel-time rerouting device. C2--C4 set a routing weight on every
edge every \ReroutePeriodS~s, $e = (1-\alpha)\,\tau + \alpha\,\bar\tau\,m/\bar m$, where $\tau$ is the edge length over the
time-mean edge speed of the last \TauAvgS~s (floored at free-flow time, as C1's device averages its speeds) and
$m$ the per-vehicle NOx cost of traversing the edge, and reroute each vehicle every \ReroutePeriodS~s with SUMO's travel-time
router on these weights, so that at $\alpha = 0$ they reduce to travel-time rerouting. C2 (placebo) replaces the
emission term with random noise matched to C4's route-change rate ($\eta = \EtaStar$); C3 measures $m$ from the
last minute ($\alpha^* = \AlphaStarReactive$); C4 forecasts it for the next minute from traffic-state features only
($\alpha^* = \AlphaStarML$). C5 replicates the legacy rule. $\alpha^*$ was chosen on validation seeds as the value
minimising NOx with mean travel time at most \AlphaTTBudgetPct\% above C1's.

\subsection{Forecaster}
Gradient-boosted trees forecast each edge's NOx per vehicle-second and mean speed for the next minute from
sensor-like features (counts, speeds, halting, waiting, their lags and neighbours, static attributes), trained on
runs with mixed controllers and tuned on validation seeds. On validation data its mean absolute error for the
per-vehicle NOx rate is \ModelMAER\ (persistence: \PersistMAER; per-edge historical mean of the training data:
\HistMAER) and for speed \ModelMAEV\ (persistence: \PersistMAEV; historical mean: \HistMAEV); hotspot PR-AUC is
\HotspotPRAUCModel\ (persistence: \HotspotPRAUCPersistence). Section~\ref{sec:forecast} evaluates it on the test
seeds.

\subsection{Statistics}
For each seed we compute the percent change of total NOx against the reference. PC2 (C4 vs C0) and PC3 (C4 vs C1)
are two-sided Wilcoxon tests; PC1 (C5 vs C0) and PC4 (C4 vs C3) are TOST equivalence tests with a margin of
$\pm$\SESOI\%. The \TestsInFamily\ tests (four comparisons times three levels) form one Holm family at \FamilywiseAlpha. The
pilot SD of paired changes (\PilotSD\%) gave $n = \ConfirmatoryN$ test seeds per condition and level, the cap of
the preregistered rule (80\% power would need $n = \NPowerSup$ for superiority and $n = \NPowerEq$ for equivalence). At
$n = \ConfirmatoryN$ the simulated power is \PowerAtNSup\% for a true change of \SESOI\% and \PowerAtNEq\% for showing
equivalence within $\pm$\SESOI\% when the true change is zero; 80\% power is reached at a true change of
\DetectableEffect\% and a margin of $\pm$\DetectableMargin\%. A non-significant test is therefore not evidence that an
effect of \SESOI\% is absent. These are planning figures based on the pilot's spread; on the test seeds the SDs of
the paired changes were \ObservedSDMin\% to \ObservedSDMax\%, so the realised precision was higher. Everything
reported beyond PC1--PC4 is exploratory, and the analyses of Sections~\ref{sec:rq5} and~\ref{sec:forecast} were added
after the results were known, following an adversarial review (\texttt{docs/DEVIATIONS.md}, entries 17--20).

\section{Results}\label{sec:results}
\subsection{The legacy result}
In the legacy notebook the rerouting runs covered \LegacyRowRatioRF\% as many vehicle-steps as the baseline, and
\LegacyRepeatUnchanged\ of \LegacyRepeatCalls\ repeated reroute calls returned the unchanged route. Its predictor
is near-deterministic: $R^2 = \LegacyRFRTwo$ on a random split and \LegacyRFRunSplitRTwo\ when split by run, because
NOx is predicted from other outputs of the same emission model at the same instant. Replicating the rule under the
matched protocol on the \LegacyTestN\ test seeds of each level, the legacy metric reports drops of
\LegacyTestLightLegacyMetric\%, \LegacyTestModerateLegacyMetric\% and \LegacyTestHeavyLegacyMetric\% (light, moderate,
heavy), while the matched metric changes by \LegacyTestLightMatchedChange\%, \LegacyTestModerateMatchedChange\% and
\LegacyTestHeavyMatchedChange\% (negative is lower NOx). The replicated rule is close to a no-op: of
\LegacyRecountCalls\ reroute calls it made on the pilot seeds, \LegacyRecountChanges\ changed a vehicle's path
(\LegacyRecountRate\%).

\subsection{Preregistered comparisons}
\PCOneLightSentence{} \PCOneModerateSentence{} \PCOneHeavySentence{}

\PCTwoLightSentence{} \PCTwoModerateSentence{} \PCTwoHeavySentence{}

\PCThreeLightSentence{} \PCThreeModerateSentence{} \PCThreeHeavySentence{}

\PCFourLightSentence{} \PCFourModerateSentence{} \PCFourHeavySentence{}

{\sloppy\PCFourDirection{}\par}

\begin{figure}[t]
\centering
\includegraphics[width=0.9\linewidth]{figures/forest_confirmatory.png}
\caption{Preregistered comparisons: mean change in total NOx with 95\% BCa intervals; dashed lines at
$\pm$\SESOI\%.}
\end{figure}

\subsection{Cost and controls (exploratory)}
Against travel-time rerouting, forecast eco-rerouting changed mean travel time by \EcoTTvsTTLight\%,
\EcoTTvsTTModerate\% and \EcoTTvsTTHeavy\% (light, moderate, heavy) and the 95th-percentile travel time by
\EcoPTTvsTTLight\%, \EcoPTTvsTTModerate\% and \EcoPTTvsTTHeavy\%. Mean departure delay, which travel time excludes,
was \DepartDelayEcoMlLight, \DepartDelayEcoMlModerate\ and \DepartDelayEcoMlHeavy~s under eco-rerouting against
\DepartDelayTtAdaptiveLight, \DepartDelayTtAdaptiveModerate\ and \DepartDelayTtAdaptiveHeavy~s under travel-time
rerouting. Total vehicle-kilometres changed by \EcoVkmVsStaticLight\%, \EcoVkmVsStaticModerate\% and
\EcoVkmVsStaticHeavy\% against no rerouting, and CO$_2$ by \EcoCOTwoVsStaticLight\%, \EcoCOTwoVsStaticModerate\% and
\EcoCOTwoVsStaticHeavy\%. Against the random-rerouting placebo, whose route-change rate was matched on validation
seeds, NOx changed by \EcoNOxVsPlaceboLight\%, \EcoNOxVsPlaceboModerate\% and \EcoNOxVsPlaceboHeavy\%, so rerouting
as such does not explain the reduction. Measured-emission rerouting changed NOx by \ReactiveNOxVsStaticLight\%,
\ReactiveNOxVsStaticModerate\% and \ReactiveNOxVsStaticHeavy\% against no rerouting, and travel-time rerouting
changed it by \TTNOxVsStaticLight\%, \TTNOxVsStaticModerate\% and \TTNOxVsStaticHeavy\%.

\subsection{Where the change occurs (exploratory, RQ5)}\label{sec:rq5}
Per-edge NOx (junction lanes excluded) shows where the reduction comes from. Against no rerouting, forecast
eco-rerouting changed NOx on residential and service streets by \ResidNOxLight\% (95\% BCa \ResidNOxCILight),
\ResidNOxModerate\% (\ResidNOxCIModerate) and \ResidNOxHeavy\% (\ResidNOxCIHeavy), and the vehicle-hours spent there by
\ResidVehHLight\%, \ResidVehHModerate\% and \ResidVehHHeavy\%; their share of edge NOx rose from
\ResidShareStaticLight, \ResidShareStaticModerate\ and \ResidShareStaticHeavy\% to \ResidShareEcoMlLight,
\ResidShareEcoMlModerate\ and \ResidShareEcoMlHeavy\%. NOx on tertiary roads changed by \TertiaryNOxLight\%,
\TertiaryNOxModerate\% and \TertiaryNOxHeavy\%, on primary and secondary roads by \PrimSecNOxLight\%,
\PrimSecNOxModerate\% and \PrimSecNOxHeavy\%. Holding each road class's NOx per vehicle-hour at the paired static
run's value, the change in where vehicles spend their time accounts for \CompositionLight\% (\CompositionCILight),
\CompositionModerate\% (\CompositionCIModerate) and \CompositionHeavy\% (\CompositionCIHeavy) of the edge-NOx
change (percentile bootstrap over seeds); against travel-time rerouting, for \CompositionVsTTLight\%,
\CompositionVsTTModerate\% and \CompositionVsTTHeavy\%. A share above 100 means NOx per vehicle-hour within road
classes rose. Measured-emission rerouting shifts NOx in the same way (Figure~\ref{fig:roadclass}).

\begin{figure}[t]
\centering
\includegraphics[width=\linewidth]{figures/road_class_shares.png}
\caption{Share of edge NOx by road class, mean over test seeds, per condition and demand level.}
\label{fig:roadclass}
\end{figure}

\begin{figure}[t]
\centering
\includegraphics[width=0.7\linewidth]{figures/map_eco_ml_minus_static_heavy.png}
\caption{Mean per-edge NOx difference, forecast eco-rerouting minus no rerouting, heavy demand (g per lane-km).}
\end{figure}

\subsection{Forecast skill on the test seeds (exploratory, RQ2)}\label{sec:forecast}
On the edge windows of the \TestRuns\ test-seed runs that recorded them (no rerouting and forecast eco-rerouting),
the forecaster's mean absolute error for NOx per vehicle-second is \TestModelMAER, against \TestHistMAER\ for the
per-edge historical mean and \TestPersistMAER\ for persistence (skill relative to the historical mean
\TestSkillHistR, 95\% cluster-bootstrap interval over runs). For speed it is \TestModelMAEV\ against \TestHistMAEV\
and \TestPersistMAEV\ (skill \TestSkillHistV), and for NOx density \TestModelMAED\ against \TestHistMAED\ and
\TestPersistMAED\ (skill \TestSkillHistD).

\subsection{Sensitivity to street coding (exploratory)}\label{sec:sensitivity}
Because the reduction depends on detours through residential streets, we repeated no rerouting, travel-time
rerouting and measured-emission rerouting on a variant of the network coded the way New York City regulates its
streets, in an experiment specified and tagged before its runs (\texttt{docs/SENSITIVITY\_PLAN.md}). The
\SensEdgesSpeedChanged\ edges that carried a road-type default speed get the city's default limit of 25~mph, which
applies to streets without a posted limit \citep{nycdot2014speed}; the \SensEdgesTagged\ edges with an OpenStreetMap
speed tag keep it; and the \SensJunctionsStop\ unsignalised junctions that a residential, unclassified or service
street enters get a stop on their minor approaches. Edges, lanes and signals are unchanged. Demand was recalibrated on
pilot seeds to the same congestion targets. Light and moderate demand reached them; heavy demand did not, because
beyond \SensDensityHeavy\ vehicles per hour per km the stop-controlled approaches jam, so heavy demand on this network
is less congested than on the original and comparisons at that level mix coding and congestion. Forecast
eco-rerouting was not run, because its forecaster takes the speed limits as inputs. The controller keeps
$\alpha^* = \AlphaStarReactive$; the same \ConfirmatoryN\ test seeds are used, with demand generated for this network.

On this network measured-emission rerouting no longer lowers NOx against travel-time rerouting. Total NOx changed
by \SensEcoReactiveVsTtAdaptiveTotalNoxGLightPosted\% (95\% BCa \SensEcoReactiveVsTtAdaptiveTotalNoxGLightPostedCI),
\SensEcoReactiveVsTtAdaptiveTotalNoxGModeratePosted\% (\SensEcoReactiveVsTtAdaptiveTotalNoxGModeratePostedCI) and
\SensEcoReactiveVsTtAdaptiveTotalNoxGHeavyPosted\% (\SensEcoReactiveVsTtAdaptiveTotalNoxGHeavyPostedCI) at light,
moderate and heavy demand, against \SensEcoReactiveVsTtAdaptiveTotalNoxGLightOriginal\%,
\SensEcoReactiveVsTtAdaptiveTotalNoxGModerateOriginal\% and \SensEcoReactiveVsTtAdaptiveTotalNoxGHeavyOriginal\% on
the original network (posted-limit minus original: \SensEcoReactiveVsTtAdaptiveTotalNoxGLightDiff,
\SensEcoReactiveVsTtAdaptiveTotalNoxGModerateDiff\ and \SensEcoReactiveVsTtAdaptiveTotalNoxGHeavyDiff\ points, intervals
\SensEcoReactiveVsTtAdaptiveTotalNoxGLightDiffCI, \SensEcoReactiveVsTtAdaptiveTotalNoxGModerateDiffCI\ and
\SensEcoReactiveVsTtAdaptiveTotalNoxGHeavyDiffCI). Against no rerouting it changed NOx by
\SensEcoReactiveVsStaticTotalNoxGLightPosted\%, \SensEcoReactiveVsStaticTotalNoxGModeratePosted\% and
\SensEcoReactiveVsStaticTotalNoxGHeavyPosted\% (original: \SensEcoReactiveVsStaticTotalNoxGLightOriginal\%,
\SensEcoReactiveVsStaticTotalNoxGModerateOriginal\%, \SensEcoReactiveVsStaticTotalNoxGHeavyOriginal\%), while
travel-time rerouting changed it by \SensTtAdaptiveVsStaticTotalNoxGLightPosted\%,
\SensTtAdaptiveVsStaticTotalNoxGModeratePosted\% and \SensTtAdaptiveVsStaticTotalNoxGHeavyPosted\%. The shift onto
residential streets disappears: with the default limit on every urban street, free-flow routes already use them
(Figure~\ref{fig:roadclassposted}), and measured-emission rerouting changed their NOx by
\SensPostedEcoReactiveVsStaticNoxResidentialServicePctLight\%, \SensPostedEcoReactiveVsStaticNoxResidentialServicePctModerate\%
and \SensPostedEcoReactiveVsStaticNoxResidentialServicePctHeavy\% against no rerouting (original:
\SensOriginalEcoReactiveVsStaticNoxResidentialServicePctLight\%,
\SensOriginalEcoReactiveVsStaticNoxResidentialServicePctModerate\% and
\SensOriginalEcoReactiveVsStaticNoxResidentialServicePctHeavy\%). Its travel-time cost against travel-time rerouting
grew to \SensEcoReactiveVsTtAdaptiveMeanTtSLightPosted\%, \SensEcoReactiveVsTtAdaptiveMeanTtSModeratePosted\% and
\SensEcoReactiveVsTtAdaptiveMeanTtSHeavyPosted\% in the mean and \SensEcoReactiveVsTtAdaptivePNineFiveTtSLightPosted\%,
\SensEcoReactiveVsTtAdaptivePNineFiveTtSModeratePosted\% and \SensEcoReactiveVsTtAdaptivePNineFiveTtSHeavyPosted\% at
the 95th percentile. Of the 540 runs, \SensInvalidRuns\ failed the validity rule; their pairs are excluded. By the reading fixed
in the plan before the runs, the reduction does not persist at any demand level; light and moderate demand, where
the two networks are equally congested, show it as clearly as heavy demand.

\begin{figure}[t]
\centering
\includegraphics[width=\linewidth]{figures/road_class_shares_posted.png}
\caption{Share of edge NOx by road class on the network with New York City's default speed limit and minor-road stop
control, mean over test seeds (compare Figure~\ref{fig:roadclass}).}
\label{fig:roadclassposted}
\end{figure}

\section{Discussion}\label{sec:discussion}
\paragraph{The original result.} The large reductions of the 2024 project are explained by its measurement: a full
baseline run was compared with rerouting runs a fraction of its length. Under a matched protocol the same rule is
equivalent to no rerouting at every demand level (PC1), and its reroute calls almost never change a path. Our own
first count of those changes was itself an artifact (a call made while a vehicle crossed a junction looked like a
change); checking that an intervention does something applies to the checking code too.

\paragraph{What eco-rerouting does in this model.} Emission-aware rerouting lowered simulated network NOx against
no rerouting, against travel-time rerouting and against a placebo matched in route-change rate, at a cost of a few
percent in travel time that is larger in the tail than in the mean. The reduction comes \CompositionWord\ from where
vehicles drive: the router moves vehicle-hours from tertiary, primary and secondary roads onto residential and
service streets, where SUMO's emission model produces less NOx per vehicle-hour at the speeds and junction
controls this network encodes. Measured-emission rerouting does the same. Two consequences follow. First, the
reduction depends on network coding. In the OpenStreetMap-derived network the arterials carry SUMO's type-default
speeds and traffic signals while residential streets have no stop control, which makes residential detours cheap.
Coded instead with New York City's default limit and stops at minor-street junctions (Section~\ref{sec:sensitivity}),
the same measured-emission controller emitted more NOx than travel-time rerouting at every demand level, and the
shift onto residential streets vanished. Forecast-driven rerouting was not run on that network (its forecaster
takes speed limits as inputs), so its result there is unknown; it optimises the same cost as the measured-emission
controller. The PC2 and PC3 reductions are therefore findings about this simulated network, not about eco-rerouting
in Flushing. Second, on the original network the emissions
that remain are moved onto streets where people live. This study has no dispersion or population model, so it says
nothing about exposure, but a routing objective that only counts total network NOx rewards exactly this
displacement.

\paragraph{The forecast.} The controller's emission cost uses the forecast NOx per vehicle-second, and on the test
seeds the forecaster \ForecastVerdictR\ a per-edge historical average of that quantity. Forecast-driven rerouting
emitted less than measured-emission rerouting by more than the equivalence margin at some demand levels, but no
preregistered condition isolates the forecast, and a learned per-edge prior is one plausible source of that
difference. We therefore do not attribute it to forecasting.

\paragraph{What carries over.} The methodological points do: matched horizons, interventions verified to act,
placebo and baseline controllers, preregistration of the comparisons and analysis code, an adversarial review
before reporting, and a check of the modelling choices a result rests on. Here the review turned a headline
reduction into a reduction with a mechanism and a side effect, and a sensitivity experiment specified before its runs
showed that the mechanism depends on how the network is coded.

\section{Limitations}
Demand is synthetic and not calibrated to New York counts; emission factors are European (HBEFA); there is no
dispersion, so emissions are not concentrations or exposure; the network and signals come from OpenStreetMap, with
SUMO's type-default speeds where OpenStreetMap has no speed tag and no stop-controlled junctions; drivers comply
fully; demand lasts one hour; the simulator is not validated against field measurements. The preregistered results
hold for the network as built; with New York City's default limit and minor-street stops the measured-emission
controller's advantage reverses (Section~\ref{sec:sensitivity}), forecast-driven rerouting was not tested on that
network, and the stop control there is an assumption, not the city's sign inventory. Road-class results cover edges
only, not junction lanes. Departure delay is reported apart from travel time. No claim here concerns health or air
quality.

\section{Reproducibility}
Code, configurations, frozen inputs (with SHA-256 hashes), the preregistration and every result file are in the
repository; each stage is one command (\texttt{python -m nycemis ...}), and every number in this paper is
generated from results files.

\section*{Acknowledgments}
The original project was co-authored by Mohammed Y.\ Hossain and Billy M.\ Kyaw and supervised by Prof.\ X.\ Huang
(NYIT). Map data \copyright\ OpenStreetMap contributors, ODbL.

\bibliographystyle{plainnat}
\bibliography{refs}
\end{document}
